\documentclass[10pt,a4paper]{amsart}
\usepackage{amsmath,amssymb,amsthm,graphicx}
\usepackage{verbatim}
\newenvironment{reference}{\comment}{\endcomment}
%\newenvironment{reference}{}{}
\newenvironment{slogan}{\comment}{\endcomment}
\newenvironment{history}{\comment}{\endcomment}

% For commutative diagrams we use Xy-pic
\usepackage[all]{xy}

% We use 2cell for 2-commutative diagrams.
\xyoption{2cell}
% Theorem environments.
%
\theoremstyle{plain}
\newtheorem{theorem}[subsection]{Theorem}
\newtheorem{proposition}[subsection]{Proposition}
\newtheorem{lemma}[subsection]{Lemma}

\theoremstyle{definition}
\newtheorem{definition}[subsection]{Definition}
\newtheorem{example}[subsection]{Example}
\newtheorem{exercise}[subsection]{Exercise}
\newtheorem{situation}[subsection]{Situation}

\theoremstyle{remark}
\newtheorem{remark}[subsection]{Remark}
\newtheorem{remarks}[subsection]{Remarks}

\numberwithin{equation}{subsection}

% Macros
%
\def\lim{\mathop{\mathrm{lim}}\nolimits}
\def\colim{\mathop{\mathrm{colim}}\nolimits}
\def\Spec{\mathop{\mathrm{Spec}}}
\def\Hom{\mathop{\mathrm{Hom}}\nolimits}
\def\Ext{\mathop{\mathrm{Ext}}\nolimits}
\def\SheafHom{\mathop{\mathcal{H}\!\mathit{om}}\nolimits}
\def\SheafExt{\mathop{\mathcal{E}\!\mathit{xt}}\nolimits}
\def\Sch{\mathit{Sch}}
\def\Mor{\mathop{\mathrm{Mor}}\nolimits}
\def\Ob{\mathop{\mathrm{Ob}}\nolimits}
\def\Sh{\mathop{\mathit{Sh}}\nolimits}
\def\NL{\mathop{N\!L}\nolimits}
\def\CH{\mathop{\mathrm{CH}}\nolimits}
\def\proetale{{pro\text{-}\acute{e}tale}}
\def\etale{{\acute{e}tale}}
\def\QCoh{\mathit{QCoh}}
\def\Ker{\mathop{\mathrm{Ker}}}
\def\Im{\mathop{\mathrm{Im}}}
\def\Coker{\mathop{\mathrm{Coker}}}
\def\Coim{\mathop{\mathrm{Coim}}}

% Boxtimes
%
\DeclareMathSymbol{\boxtimes}{\mathbin}{AMSa}{"02}

%
% Macros for moduli stacks/spaces
%
\def\QCohstack{\mathcal{QC}\!\mathit{oh}}
\def\Cohstack{\mathcal{C}\!\mathit{oh}}
\def\Spacesstack{\mathcal{S}\!\mathit{paces}}
\def\Quotfunctor{\mathrm{Quot}}
\def\Hilbfunctor{\mathrm{Hilb}}
\def\Curvesstack{\mathcal{C}\!\mathit{urves}}
\def\Polarizedstack{\mathcal{P}\!\mathit{olarized}}
\def\Complexesstack{\mathcal{C}\!\mathit{omplexes}}
% \Pic is the operator that assigns to X its picard group, usage \Pic(X)
% \Picardstack_{X/B} denotes the Picard stack of X over B
% \Picardfunctor_{X/B} denotes the Picard functor of X over B
\def\Pic{\mathop{\mathrm{Pic}}\nolimits}
\def\Picardstack{\mathcal{P}\!\mathit{ic}}
\def\Picardfunctor{\mathrm{Pic}}
\def\Deformationcategory{\mathcal{D}\!\mathit{ef}}
\usepackage{cleveref}
\crefname{originalreference}{original reference}{original references}
\Crefname{originalreference}{Original reference}{Original references}
\makeatletter
\newlabel{sites-cohomology-lemma-Leray}{{0732}{source}{}{sourcepointer-1}{}}
\newlabel{sites-cohomology-lemma-Leray@cref}{{[originalreference][0][]0732}{[1][0][]source}{}{}{}}
\makeatother
\crefname{originallemma}{original lemma}{original lemmas}
\Crefname{originallemma}{Original lemma}{Original lemmas}
\makeatletter
\newlabel{lemma-adjoint-kernel-parasitic}{{0779}{source}{}{sourcepointer-2}{}}
\newlabel{lemma-adjoint-kernel-parasitic@cref}{{[originallemma][0][]0779}{[1][0][]source}{}{}{}}
\makeatother
\crefname{originallemma}{original lemma}{original lemmas}
\Crefname{originallemma}{Original lemma}{Original lemmas}
\makeatletter
\newlabel{lemma-pushforward-parasitic}{{0775}{source}{}{sourcepointer-3}{}}
\newlabel{lemma-pushforward-parasitic@cref}{{[originallemma][0][]0775}{[1][0][]source}{}{}{}}
\makeatother
\title{Leray spectral sequence for the derived pushforward of the qc category on algebraic stacks}
\author{The Stacks Project Authors}
\date{stacks-project:stacks-cohomology.tex:local-snapshot-20261012; GNU FDL 1.2 or later}
\begin{document}
\maketitle
\noindent\textit{Editorial note.} All delivered mathematical text is the author's own, copied verbatim from the pinned Stacks Project chapter file stacks-cohomology.tex; no statement or proof step was rewritten, reordered or removed.\par\medskip
\noindent\textit{Editorial note.} The counted span cites the lemmas recorded in assembly.external\_references with canonical tags as navigation pointers.\par\medskip
\noindent\textit{Allowed standard prerequisites (editorial).}\begin{itemize}
\item Leray spectral sequence for ringed sites. Stacks Project, Cohomology on Sites, Leray spectral sequence (as specialized in sites-cohomology). Provides the two spectral sequences being compared.
\item Kernel and cokernel of the adjunction map are parasitic. Stacks Project, Cohomology of Algebraic Stacks, Lemma stacks-cohomology-lemma-adjoint-kernel-parasitic. Identifies the comparison map's kernel and cokernel.
\item Cohomology of parasitic modules vanishes. Stacks Project, Cohomology of Algebraic Stacks, Lemma stacks-cohomology-lemma-pushforward-parasitic. Ensures the comparison map induces isomorphisms on cohomology.
\end{itemize}\medskip
\ifdefined\hypertarget\hypertarget{sourcepointer-1}{}\fi
\noindent\textit{Source reference.} Original reference 0732: sources/stacks-chow/tags.txt, source lines 9046--9046. Cross-chapter standard prerequisite cited by the author inside the counted span; it is used as a pointer and its own proof is not part of this item.\par\medskip
\ifdefined\hypertarget\hypertarget{sourcepointer-2}{}\fi
\noindent\textit{Source reference.} Original lemma 0779: sources/stacks-chow/tags.txt, source lines 9194--9194. Cross-chapter standard prerequisite cited by the author inside the counted span; it is used as a pointer and its own proof is not part of this item.\par\medskip
\ifdefined\hypertarget\hypertarget{sourcepointer-3}{}\fi
\noindent\textit{Source reference.} Original lemma 0775: sources/stacks-chow/tags.txt, source lines 9190--9190. Cross-chapter standard prerequisite cited by the author inside the counted span; it is used as a pointer and its own proof is not part of this item.\par\medskip
\begin{lemma}
\label{lemma-leray}
Let $f : \mathcal{X} \to \mathcal{Y}$
be a quasi-compact and quasi-separated morphism of algebraic stacks.
Let $\mathcal{F}$ be a quasi-coherent sheaf on $\mathcal{X}$. Then
there exists a spectral sequence with $E_2$-page
$$
E_2^{p, q} = H^p(\mathcal{Y}, R^qf_{\QCoh, *}\mathcal{F})
$$
converging to $H^{p + q}(\mathcal{X}, \mathcal{F})$.
\end{lemma}

\begin{proof}
By Cohomology on Sites, Lemma \ref{sites-cohomology-lemma-Leray}
the Leray spectral sequence with
$$
E_2^{p, q} = H^p(\mathcal{Y}, R^qf_*\mathcal{F})
$$
converges to $H^{p + q}(\mathcal{X}, \mathcal{F})$. The kernel and cokernel
of the adjunction map
$$
R^qf_{\QCoh, *}\mathcal{F} \longrightarrow R^qf_*\mathcal{F}
$$
are parasitic modules on $\mathcal{Y}$
(Lemma \ref{lemma-adjoint-kernel-parasitic})
hence have vanishing cohomology
(Lemma \ref{lemma-pushforward-parasitic}).
It follows formally that
$H^p(\mathcal{Y}, R^qf_{\QCoh, *}\mathcal{F}) =
H^p(\mathcal{Y}, R^qf_*\mathcal{F})$ and we win.
\end{proof}

\end{document}
